% =====================================================================
%  Final Report — Computer Security Project   (Deliverable C)
%  Recommended length: 6 pages
%
%  HOW TO USE THIS FILE
%  1. Paste into a new Overleaf project. It compiles as-is.
%  2. Search this file for  FILLIN  (shown in red in the PDF) and replace
%     every occurrence with real numbers from your Colab notebook.
%  3. The five tables below mirror the LaTeX tables the notebook prints
%     under "--- LaTeX ---". Paste those blocks over the placeholders.
%  4. Replace \figbox{...} with \includegraphics{...} once you have
%     exported the PNGs from the notebook (keep the \caption).
% =====================================================================
\documentclass[11pt]{article}
\usepackage[margin=1in]{geometry}
\usepackage{graphicx, booktabs, xcolor, hyperref, amsmath, array, enumitem, caption}
\usepackage[T1]{fontenc}
\usepackage{lmodern}

\hypersetup{colorlinks=true, linkcolor=blue!50!black, citecolor=blue!50!black, urlcolor=blue!50!black}

% ---- edit-me markers -------------------------------------------------
\newcommand{\fillin}[1]{\textcolor{red}{\textbf{[#1]}}}
% ---- placeholder figure box (replace with \includegraphics) ----------
\newcommand{\figbox}[2]{%
  \begin{center}
    \fbox{\parbox[c][0.30\textheight][c]{0.96\textwidth}{\centering\large\textit{#1}}}
    \\[4pt]\textbf{Figure placeholder:} #2
  \end{center}}

\title{\textbf{Can Abnormal IoT Traffic Be Detected?}\\[4pt]
\large Behavioral Detection of Botnet-Infected IoT Devices using Machine Learning}
\author{\fillin{Your Name}\\ \fillin{Course Name / Number}\\ \fillin{Instructor}}
\date{\fillin{Date}}

\begin{document}
\maketitle

% =====================================================================
\section{Introduction}
% =====================================================================
The Internet of Things has expanded far faster than the security controls
available to protect it. Tens of billions of devices --- cameras, doorbells,
thermostats, baby monitors, industrial sensors --- are now connected to home
and enterprise networks, and they share three properties that make them
attractive targets: they ship with default or hard-coded credentials, they are
rarely patched, and they cannot run endpoint security software. Once
compromised they are routinely conscripted into botnets such as \textbf{Mirai}
and \textbf{BASHLITE (Gafgyt)}, which have been used to launch some of the
largest distributed denial-of-service (DDoS) attacks on record
\cite{kolias2017ddos}.

Because host-based defenses are not available for these devices, the only
control that scales is \emph{network-side monitoring}: observing a device's
traffic from a gateway or intrusion detection system (IDS) and flagging
behaviour that deviates from its own established baseline. This project asks
whether that is achievable in practice from behavioural traffic features
alone.

\medskip
\noindent\textbf{Research question:}
\emph{Can behavioral network-traffic features accurately distinguish benign
IoT traffic from botnet-infected IoT traffic?}

The analysis uses the \textbf{N-BaIoT} dataset \cite{kashif2020nbaiot},
which contains real traffic from nine commercial IoT devices captured before
and after authentic infection with Mirai and BASHLITE. Four machine learning
methods are compared --- Logistic Regression, Random Forest, XGBoost, and an
Isolation Forest anomaly detector --- and the results are interpreted not
merely as accuracy figures but in terms of what they mean for a security
operator: which attacks are caught, which slip through, and what the false
positive and false negative rates cost in the real world.

% =====================================================================
\section{Security Problem}
% =====================================================================
The specific problem investigated is the \textbf{detection of compromised IoT
devices that have joined a botnet, using only their network traffic
behaviour}.

Mirai and BASHLITE follow a well-documented lifecycle: \emph{propagation}
(scanning for new victims), \emph{infection} (brute-forcing default Telnet
credentials), \emph{command and control}, and \emph{execution} (launching
floods on command). This project targets the \emph{execution} stage, because
it is the stage that causes external harm: at that moment the device is
actively attacking someone else's network from inside yours.

Three constraints shape the solution:

\begin{enumerate}[itemsep=2pt]
  \item \textbf{No endpoint agent.} IoT devices lack the storage, memory and
        power budget for host-based detection, and users rarely patch them.
  \item \textbf{Heterogeneity.} A single site may host dozens of device
        models, each with its own traffic fingerprint.
  \item \textbf{Scale.} Manual inspection of traffic is impossible; detection
        must be automated and cheap enough to run continuously at the edge.
\end{enumerate}

A behavioural, network-based detector satisfies all three: it runs on the
access point or gateway, consumes no device resources, and needs no vendor
cooperation.

\medskip
A critical point frames everything that follows. For a security tool,
\textbf{accuracy is not the objective}. A false negative leaves a compromised
device attacking from inside the network; a false positive can take a
legitimate --- and possibly safety-critical --- device offline and contribute
to analyst alert fatigue. Both error types are therefore reported explicitly
throughout, and their asymmetric consequences are analysed in
Section~\ref{sec:analysis}.

% =====================================================================
\section{Dataset Description}
% =====================================================================
The dataset used is \textbf{N-BaIoT Dataset to Detect IoT Botnet Attacks}
\cite{kashif2020nbaiot}, available at
\url{https://www.kaggle.com/datasets/mkashifn/nbaiot-dataset}. It originates
from the Ben-Gurion University lab study described by Meidan et al.
\cite{meidan2018nbaiot} and is also mirrored at UCI \cite{uci442}.

\begin{table}[h]
\centering
\caption{Dataset summary.}
\label{tab:dataset}
\begin{tabular}{@{} p{0.32\textwidth} p{0.6\textwidth} @{}}
\toprule
\textbf{Field} & \textbf{Value} \\
\midrule
Name & N-BaIoT Dataset to Detect IoT Botnet Attacks \\
Kaggle URL & \url{https://www.kaggle.com/datasets/mkashifn/nbaiot-dataset} \\
Records & 7,062,606 on disk; 150,304 sampled \\
Features & 115 behavioural statistics \\
Target variable & \texttt{label}: \texttt{0 = benign}, \texttt{1 = malicious} \\
Devices & 9 commercial IoT devices \\
Attack types & Mirai and BASHLITE variants: UDP, TCP, ACK, SYN, UDPplain, Scan, Junk \\
\bottomrule
\end{tabular}
\end{table}

Rather than raw packets, each record is a \emph{behavioral snapshot}: 23
statistics computed over each of five time windows (100\,ms, 500\,ms, 1.5\,s,
10\,s, 1\,min), giving $23 \times 5 = 115$ features. The statistics are
aggregated by source IP, source MAC-IP, channel, and socket.

\begin{table}[h]
\centering
\caption{Description of important features.}
\label{tab:features}
\begin{tabular}{@{} p{0.30\textwidth} p{0.30\textwidth} p{0.34\textwidth} @{}}
\toprule
\textbf{Feature group} & \textbf{Statistics} & \textbf{Security relevance} \\
\midrule
Packet size (outbound) & mean, variance & Floods replace normal payloads with
uniform junk or oversized packets \\
Packet count & count & Scanning and flooding push packet rates far above the
device baseline \\
Packet jitter & mean, variance, number & Flood generators are mechanically
regular; command-and-control beacons alter inter-arrival timing \\
Packet size (in + out) & magnitude, radius, covariance, correlation &
Captures the shift from bidirectional request/response to one-way attack
traffic \\
\bottomrule
\end{tabular}
\end{table}

\subsection{Data Exploration}\label{sec:eda}
Exploratory analysis covered class distribution, per-device composition,
missing values, duplicate records, feature distributions, correlation
structure, and per-feature effect sizes.

\figbox{Class distribution + records per device (bar charts)}
{Generated by the notebook's \texttt{Class distribution + device composition} cell.}

\figbox{Correlation heatmap of the highest-variance features}
{Generated by the notebook's \texttt{Correlation analysis} cell.}

\noindent Key observations:

\begin{itemize}[itemsep=2pt]
  \item \textbf{Class distribution.} The 150,304-row sample holds
        \textbf{12,113 benign (8.1\%)} against \textbf{138,191 malicious
        (91.9\%)} --- a ratio of \textbf{11.4:1}. A production network is the
        exact opposite, so this inverted prior is the single largest caveat in
        the study and is revisited in Section~\ref{sec:limitations}.
  \item \textbf{Missing values and duplicates.} \textbf{Zero missing values}
        across all 115 features, and \textbf{96 exact duplicate records}
        (0.06\%: 35 benign, 61 malicious).
        Duplicates arise because flood generators emit byte-identical flow
        snapshots repeatedly; they were removed before splitting.
  \item \textbf{Distribution shape.} Flow statistics are strongly
        right-skewed, motivating log-scale plots and favouring tree-based
        models over linear ones.
  \item \textbf{Correlation.} Large blocks of near-duplicate features exist
        because the same 23 statistics repeat across five time windows;
        \textbf{334} feature pairs exceeded $|r| > 0.95$ (mean absolute
        correlation 0.272) --- a direct consequence of repeating the same 23
        statistics across five windows.
  \item \textbf{Separating features.} The largest effect sizes are all
        \texttt{weight} statistics --- the number of packets observed in the
        recent window. \texttt{MI\_dir\_L0.1\_weight} averages \textbf{14.5}
        on benign traffic versus \textbf{3,689.6} when infected
        (\textbf{254$\times$, $d = 2.56$}); \texttt{MI\_dir\_L0.01\_weight}
        averages 68.9 versus 21,472.4 ($311\times$). An infected device simply
        moves two to three orders of magnitude more packets in the shortest
        windows.
\end{itemize}

% =====================================================================
\section{Data Preprocessing}
% =====================================================================
Each step below is necessary for a specific reason; all data-dependent
transforms are fitted on the training partition only.

\begin{table}[h]
\centering
\caption{Preprocessing steps and their justification.}
\label{tab:preprocessing}
\begin{tabular}{@{} p{0.28\textwidth} p{0.64\textwidth} @{}}
\toprule
\textbf{Step} & \textbf{Why it is necessary} \\
\midrule
Remove duplicates & Byte-identical flood snapshots would otherwise appear in
both train and test, inflating accuracy and wasting memory. Performed
\emph{before} the split so no duplicate can leak across partitions. \\
\addlinespace
Train/test split & A stratified $80/20$ split keyed on
$\text{device} \times \text{label}$ guarantees every device and both classes
appear in both partitions, making the per-device results in
Section~\ref{sec:results} valid. \\
\addlinespace
Missing values & Median imputation, robust to the heavy right skew of flow
statistics. Fitted on training data only. \\
\addlinespace
Feature selection & Near-zero-variance and highly correlated features
($|r| > 0.95$) slow training, destabilise logistic-regression coefficients,
and dilute importance rankings. Reduces \textbf{115} to \textbf{36}
features. \\
\addlinespace
Encoding categorical & The device identity is a legitimate but \emph{nominal}
feature; one-hot encoding avoids implying a non-existent ordering. \\
\addlinespace
Feature scaling & Standardisation puts all features on one range so
gradient-based models converge and no single large-magnitude statistic
dominates. Fitted on training data only. \\
\addlinespace
Imbalanced classes & \texttt{class\_weight='balanced'} penalises errors on the
minority class proportionally to its rarity; recall, F1, FPR and FNR are
reported instead of accuracy alone. \\
\bottomrule
\end{tabular}
\end{table}

% =====================================================================
\section{Methodology}
% =====================================================================
Four methods were selected to span complementary assumptions about what
\emph{abnormal} means, exceeding the two-method minimum:

\begin{table}[h]
\centering
\caption{Methods and how each addresses the security problem.}
\label{tab:methods}
\begin{tabular}{@{} p{0.20\textwidth} p{0.16\textwidth} p{0.56\textwidth} @{}}
\toprule
\textbf{Method} & \textbf{Type} & \textbf{Role in the analysis} \\
\midrule
Logistic Regression & Supervised, linear & Interpretable baseline. Each
coefficient is a signed weight on a traffic statistic, so the direction and
magnitude of each behaviour's contribution can be read directly. \\
\addlinespace
Random Forest & Supervised, ensemble & Captures non-linear and interaction
effects (jitter matters only when packet rate is also elevated) and supplies a
built-in feature-importance ranking. \\
\addlinespace
XGBoost & Supervised, boosting & Strongest tabular learner; provides
calibrated scores and SHAP values for explainability. \\
\addlinespace
Isolation Forest & \textbf{Unsupervised anomaly detection} & Trained on
\emph{benign traffic only} --- it never sees an attack. It flags statistically
unusual flows, so it can in principle detect novel behaviour that no labelled
training set contains. Closest analogue to a real-world IDS. \\
\bottomrule
\end{tabular}
\end{table}

All models were implemented in Python using scikit-learn, XGBoost and SHAP
\cite{pedregosa2011sklearn, chen2017xgboost, liu2008isolation,
lundberg2017shap, breiman2001random}. Every random seed was fixed at
\texttt{SEED = 42} for reproducibility.

% =====================================================================
\section{Experimental Results}\label{sec:results}
% =====================================================================
After de-duplication the corpus held 150,208 rows, split 80/20 into
\textbf{120,166 training} and \textbf{30,042 test} rows --- the test partition
containing only \textbf{2,416 benign} and \textbf{27,626 malicious} flows. That
benign count is small enough to matter: see Section~\ref{sec:limitations}.

\begin{table}[h]
\centering
\caption{Overall model performance on the held-out test set.}
\label{tab:overall}
\begin{tabular}{@{} l c c c c c c c @{}}
\toprule
\textbf{Model} & \textbf{Acc.} & \textbf{Prec.} & \textbf{Recall} &
\textbf{F1} & \textbf{AUC} & \textbf{FPR} & \textbf{FNR} \\
\midrule
Logistic Regression & 0.9992 & 0.9996 & 0.9996 & 0.9996 & 0.9989 & 0.0046 & 0.0004 \\
Random Forest       & 0.9999 & 1.0000 & 0.9999 & 1.0000 & 1.0000 & 0.0000 & 0.0001 \\
XGBoost             & 1.0000 & 1.0000 & 1.0000 & 1.0000 & 1.0000 & 0.0000 & 0.0000 \\
Isolation Forest    & 0.2978 & 0.9706 & 0.2438 & 0.3897 & 0.8835 & 0.0844 & 0.7562 \\
\bottomrule
\end{tabular}
\end{table}

\figbox{Confusion matrices (one per model)}
{Notebook cell: \texttt{Confusion matrices}.}

\figbox{ROC curves with AUC values}
{Notebook cell: \texttt{ROC curves}.}

\subsection*{Detection rate by attack type}
\begin{table}[h]
\centering
\caption{Detection rate by attack type. Attack rows = recall; the benign row = specificity.}
\label{tab:attacks}
\begin{tabular}{@{} l c c c @{}}
\toprule
\textbf{Attack} & \textbf{Random Forest} & \textbf{XGBoost} & \textbf{Isolation Forest} \\
\midrule
mirai\_udpplain & 1.0000 & 1.0000 & 0.6523 \\
mirai\_udp      & 1.0000 & 1.0000 & 0.6514 \\
mirai\_ack      & 1.0000 & 1.0000 & 0.6491 \\
gafgyt\_junk    & 1.0000 & 1.0000 & \textbf{0.0274} \\
gafgyt\_scan    & 1.0000 & 1.0000 & \textbf{0.0226} \\
gafgyt\_combo   & 1.0000 & 1.0000 & \textbf{0.0024} \\
gafgyt\_udp     & 0.9998 & 0.9998 & \textbf{0.0002} \\
gafgyt\_tcp     & 0.9997 & 1.0000 & \textbf{0.0000} \\
mirai\_scan     & 1.0000 & 1.0000 & \textbf{0.0000} \\
mirai\_syn      & 1.0000 & 1.0000 & \textbf{0.0000} \\
benign (specificity) & 1.0000 & 1.0000 & 0.9156 \\
\bottomrule
\end{tabular}
\end{table}

\subsection*{Operational translation of the error rates}
\begin{table}[h]
\centering
\caption{Operational impact at 1{,}000{,}000 benign flows per day.}
\label{tab:ops}
\begin{tabular}{@{} l c c c @{}}
\toprule
\textbf{Model} & \textbf{FPR} & \textbf{False alerts / day} & \textbf{Missed attacks / day} \\
\midrule
Logistic Regression & 0.0046 & 4{,}553 & 4.3 \\
Random Forest       & 0.0000 & 0 & 0.7 \\
XGBoost             & 0.0000 & 0 & 0.4 \\
Isolation Forest    & 0.0844 & 84{,}437 & 7{,}562 \\
\bottomrule
\end{tabular}
\end{table}

\noindent This table is the bridge between benchmark metrics and deployment
relevance: an FPR that looks negligible in a paper becomes \textbf{4,553
false alerts per day} for Logistic Regression and \textbf{84,437 per day} for
the Isolation Forest. Even the measured ``zero'' of Random Forest and XGBoost
only bounds the true rate to $\lesssim 0.124\%$ at 95\% confidence (rule of
three over 2,416 benign flows), i.e.\ up to $\sim$1,242 alerts per day. This is
what determines whether an operator keeps the system switched on.

\subsection*{Per-device performance}
\begin{table}[h]
\centering
\caption{F1-score per device (pooled model evaluated on each device separately).}
\label{tab:per_device}
\begin{tabular}{@{} l c c c c @{}}
\toprule
\textbf{Device} & \textbf{Log. Reg.} & \textbf{Rand. Forest} & \textbf{XGBoost} & \textbf{Isol. Forest} \\
\midrule
Danmini Doorbell                    & 0.9996 & 1.0000 & 1.0000 & 0.4178 \\
Ecobee Thermostat                   & 0.9996 & 0.9999 & 0.9999 & 0.4874 \\
Ennio Doorbell                      & 1.0000 & 1.0000 & 1.0000 & \textbf{0.0015} \\
Philips B120N10 Baby Monitor        & 0.9999 & 1.0000 & 1.0000 & 0.4216 \\
Provision PT 737E Security Camera   & 0.9997 & 1.0000 & 1.0000 & 0.2785 \\
Provision PT 838 Security Camera    & 0.9994 & 0.9998 & 1.0000 & 0.2719 \\
Samsung SNH 1011 N Webcam           & 0.9982 & 1.0000 & 1.0000 & \textbf{0.0015} \\
SimpleHome XCS7 1002 WHT Cam.       & 0.9994 & 1.0000 & 1.0000 & 0.4279 \\
SimpleHome XCS7 1003 WHT Cam.       & 0.9999 & 1.0000 & 1.0000 & 0.5894 \\
\bottomrule
\end{tabular}
\end{table}

\noindent The pooled supervised model generalises across the fleet: XGBoost
scores 1.0000 on eight of nine devices (0.9999 on the Ecobee), Random Forest
stays at or above 0.9998 everywhere, and Logistic Regression only dips to
0.9982 on the Samsung webcam. The Isolation Forest, by contrast, is worse than
random on two devices (\textbf{F1 = 0.0015} on both the Ennio Doorbell and the
Samsung webcam) and never exceeds 0.59. Two of the nine devices --- the Ennio
Doorbell and the Samsung webcam, which ship only six attack files each
(BASHLITE only) --- collapse entirely, a strong argument that a single pooled
anomaly threshold cannot serve a heterogeneous fleet.
\figbox{F1, FPR and FNR broken down by device}
{Notebook cell: \texttt{Per-device performance}.}

% =====================================================================
\section{Security Analysis}\label{sec:analysis}
% =====================================================================
\subsection{What the detector actually detects}
The models distinguish benign traffic from traffic produced by devices running
Mirai or BASHLITE --- concretely, a device that is scanning for new victims,
flooding a target, or transmitting junk data on command. Detecting this stage
matters because it is the stage that causes external harm.

They do it without ever inspecting a payload. The largest effect sizes in
Section~\ref{sec:eda} are all \emph{packet-volume} statistics:
\texttt{MI\_dir\_L0.1\_weight} averages \textbf{14.5} on benign traffic and
\textbf{3,689.6} when infected --- \textbf{254$\times$} more packets in a
0.1\,s window. The detector is effectively noticing that a device whose job is
to measure temperature or stream a doorbell view has begun generating traffic
at flood rates. Task-oriented devices are predictable, so the deviation is
enormous rather than subtle; that is what makes IoT botnet detection tractable
at all.

\subsection{Which attacks are difficult to detect}
The answer splits sharply by method.

\textbf{Supervised models miss essentially nothing.} Random Forest and XGBoost
reached recall 1.0000 on \emph{nine of the ten} attack types; the only
sub-perfect cells were \texttt{gafgyt\_tcp} (RF 0.9997) and \texttt{gafgyt\_udp}
(RF and XGBoost 0.9998). Logistic Regression missed 12 instances in total.

\textbf{Unsupervised detection fails on exactly the attacks that matter most.}
The Isolation Forest caught only the three highest-rate Mirai floods
(\texttt{mirai\_udpplain} 0.6523, \texttt{mirai\_udp} 0.6514,
\texttt{mirai\_ack} 0.6491) and scored:
\textbf{0.0000} on \texttt{mirai\_scan} and \texttt{mirai\_syn},
\textbf{0.0000} on \texttt{gafgyt\_tcp}, \textbf{0.0002} on
\texttt{gafgyt\_udp}, \textbf{0.0024} on \texttt{gafgyt\_combo},
\textbf{0.0226} on \texttt{gafgyt\_scan} and \textbf{0.0274} on
\texttt{gafgyt\_junk}. Overall it missed \textbf{three out of every four attack
instances} (FNR 0.7562).

Two conclusions follow. First, \textbf{scanning is both the hard case and the
dangerous one}: it is the \emph{propagation} phase, so a miss lets the botnet
keep recruiting while a flood is loud but self-limiting. Second, \textbf{``anomaly
detection'' is not automatically good at detecting attacks} --- benign IoT
traffic is itself highly variable (a camera streaming video legitimately
produces huge packet volumes), so high volume alone is not anomalous. The gap
with the original paper's 100\% TPR is explained by configuration: they used
one tuned deep autoencoder \emph{per device} with a threshold and voting window
optimised on held-out benign data, whereas this study used a single pooled
Isolation Forest with default hyperparameters.

\subsection{Which features are most important}
XGBoost gain importance is dominated by two features that together carry
\textbf{87\% of the total}: \texttt{HpHp\_L0.01\_radius} (0.6898) and
\texttt{HH\_L0.01\_pcc} (0.1829), followed by \texttt{HH\_jit\_L5\_mean}
(0.0442), \texttt{HH\_L0.01\_weight} (0.0255) and
\texttt{MI\_dir\_L0.01\_weight} (0.0210).

Their vocabulary is consistent and physically meaningful:
\textbf{\texttt{L0.01} is the shortest window (10\,ms)}, so detection is driven
by the most recent few milliseconds rather than long-run averages;
\textbf{\texttt{HpHp} is the socket level} and \textbf{\texttt{HH} the channel
level}; and \texttt{radius} is the root-sum-of-squares of the two streams'
variances --- a direct measure of \emph{burstiness on one connection}, which
attack floods exhibit and normal request/response traffic does not.
\texttt{HH\_jit} captures jitter: flood generators and C\&C beacons are
mechanically timed, real sensors are not.

SHAP on individual predictions agrees and shows both directions:
\texttt{HH\_jit\_L5\_mean} contributed \textbf{+6.75} to the example malicious
prediction (every other feature $\leq +0.54$), while
\texttt{MI\_dir\_L0.01\_weight} contributed \textbf{$-5.05}$ to the benign one.

\medskip
\noindent\textbf{Operational reading:} the decision rests on burstiness on a
single socket within the last 10\,ms, corroborated by jitter over the last 5\,s
and packet volume over the last 10\,ms --- all simple incremental statistics
cheap enough to compute continuously on gateway hardware.

\subsection{Consequences of false positives}
A false positive labels a legitimate device as compromised. \textbf{Measured:}
Logistic Regression produced 11 false positives on 2,416 benign flows (FPR
0.0046 $\rightarrow$ 4,553 alerts/day), Random Forest and XGBoost produced 0,
and the Isolation Forest produced 204 (FPR 0.0844 $\rightarrow$ 84,437
alerts/day).
\begin{itemize}[itemsep=2pt]
  \item \textbf{Alert fatigue.} Two alerts per minute (Logistic Regression) is
        already beyond human triage; 84,437 per day is not a security product,
        it is a denial of service against your own SOC.
  \item \textbf{Availability loss.} If the response is automatic quarantine, a
        false positive takes a camera, doorbell or thermostat offline. For a
        baby monitor or security camera that is a \emph{physical-safety and
        privacy outage}, not merely an IT ticket.
  \item \textbf{Why ``0'' is not really 0.} Only 2,416 benign flows reached the
        test set. A measured rate of zero over 2,416 trials bounds the true FPR
        only to $\approx 3/2416 = 0.124\%$ at 95\% confidence --- up to
        $\sim$1,242 false alerts/day. \textbf{The benign test set is too small
        to certify the FPR that deployment requires.}
\end{itemize}

\subsection{Consequences of false negatives}
A false negative means a compromised device stays on the network.
\textbf{Measured:} Logistic Regression FN $=$ 12, Random Forest FN $=$ 2,
XGBoost FN $=$ 0, and Isolation Forest FN $=$ \textbf{20,890} (FNR 0.7562) ---
about \textbf{7,562 missed attacks per day} at 10,000 attack flows/day.
\begin{itemize}[itemsep=2pt]
  \item it continues participating in \textbf{DDoS attacks} against third
        parties, creating legal and reputational exposure for the network owner;
  \item it keeps \textbf{scanning and propagating}, recruiting further devices;
  \item it provides an \textbf{inside foothold} for later lateral movement and
        data exfiltration;
  \item in a home, an infected camera is a \textbf{surveillance and privacy
        breach}.
\end{itemize}

\medskip
\noindent This asymmetry is the central security argument of the project:
\textbf{missed detection is far more damaging than a false alarm}, which is
why recall and FNR --- rather than accuracy --- are treated as the primary
measures of success.

\subsection{Use in a real security system}
The approach would be deployed at the \textbf{Wi-Fi access point or edge
gateway}, where all device traffic already transits and no agent is required on
the device --- which is the entire reason network-side detection is the right
choice for IoT: no firmware patch, no vendor cooperation, no drain on hardware
that may have kilobytes of headroom.

The per-device breakdown settles the model-topology question. A single pooled
\emph{supervised} model is sufficient: XGBoost reaches F1 = 1.0000 on eight of
nine devices and 0.9999 on the ninth. A single pooled \emph{anomaly} model is
not: the Isolation Forest collapses to F1 = 0.0015 on two devices, supporting
the original study's argument for \emph{one baseline per device} whenever
unsupervised detection is used.

A practical response policy uses \textbf{tiered thresholds} chosen from the
FPR/FNR trade-off in Table~\ref{tab:ops} rather than a defaulted 0.5 cut:
high-confidence predictions trigger automatic quarantine, medium-confidence
ones are routed to an analyst. Because firmware updates change traffic
patterns, models require periodic re-baselining against benign traffic or the
FPR will drift upward. The feature budget is also favourable --- the top ten
importances are all simple incremental statistics over 10\,ms--5\,s windows,
cheap enough for continuous scoring on commodity gateway hardware. Finally, a
supervised detector can only ever recognise labelled behaviour; a second,
per-device anomaly model is the natural complement for attacks outside the
training repertoire, provided it is tuned as carefully as Section~\ref{sec:analysis}
argues it must be.

\subsection{Limitations}
See Section~\ref{sec:limitations}.

% =====================================================================
\section{Limitations}\label{sec:limitations}
% =====================================================================
\begin{enumerate}[itemsep=3pt]
  \item \textbf{Lab versus production.} N-BaIoT was captured in an isolated
        lab with nine devices and two botnets. Real networks add more devices,
        more protocols, encrypted traffic, and far more benign variety.
  \item \textbf{Inverted class prior --- the largest caveat.} The sample is
        \textbf{91.9\% malicious} (11.4:1). A production network is the
        opposite. A precision of 1.0000 measured here would collapse under a
        realistic prior: with FPR 0.0046 and a 0.1\% attack rate, the majority
        of ``alerts'' would be benign traffic rather than attacks. Every
        precision figure in Section~\ref{sec:results} is optimistic.
  \item \textbf{Benign test set too small to certify FPR.} Only 2,416 benign
        flows were held out. Zero observed false positives bounds the true rate
        only to $\lesssim 0.124\%$ (95\% CI) --- up to $\sim$1,242 false
        alerts/day at one million flows.
  \item \textbf{Naive anomaly-detector configuration.} One pooled Isolation
        Forest, default \texttt{contamination}, no per-device baseline, no
        threshold tuning, no majority-vote window. Its poor result is a property
        of this configuration, not a verdict on anomaly detection.
  \item \textbf{Random rather than chronological split.} Rows were split
        randomly, so near-identical consecutive flow snapshots may appear on
        both sides of the split. This inflates results and is stated openly.
  \item \textbf{No adversarial testing.} An attacker aware that detection
        relies on jitter and rate statistics can shape traffic to evade them
        (rate-limiting, jitter padding). Robustness to an adaptive adversary is
        untested.
  \item \textbf{Two botnet families only.} Novel malware with different
        behaviour is what a supervised model will miss. The unsupervised
        Isolation Forest partially addresses this but is weaker on subtle
        attacks.
  \item \textbf{Content-blind features.} Behavioural statistics abstract away
        payload, so the detector can flag deviation from baseline but cannot by
        itself say whether a high-volume flow is an attack or legitimate
        streaming.
  \item \textbf{Offline training.} An operational system needs streaming
        feature extraction and incremental scoring, which the original paper
        addresses but this study does not implement.
\end{enumerate}

% =====================================================================
\section{Conclusion}
% =====================================================================
\textbf{Answer to the research question:}
Yes --- with important caveats. Behavioural network-traffic features distinguish
benign from botnet-infected IoT traffic almost perfectly on this benchmark:
\textbf{XGBoost F1 = 1.0000 and ROC-AUC = 1.0000}, Random Forest F1 = 1.0000
and Logistic Regression F1 = 0.9996, consistently across \textbf{all nine
devices} (XGBoost F1 = 1.0000 on eight and 0.9999 on the Ecobee thermostat).

The security conclusions are as important as the metric:

\begin{enumerate}[itemsep=3pt]
  \item Network-side detection is viable for devices that cannot protect
        themselves --- no agent, no patch, no user action required.
  \item \textbf{Floods are easy; scanning is the weak point --- and it is the
        one that matters.} Supervised models caught 9 of 10 attack types at
        1.0000 recall, but the Isolation Forest detected
        \texttt{mirai\_scan} and \texttt{mirai\_syn} \textbf{0\%} of the time
        and \texttt{gafgyt\_scan} at 2.3\%. Scanning is the propagation stage,
        so those are the misses that let a botnet grow.
  \item \textbf{The signal is volume and burstiness, not malware content.}
        Infected devices move 100--300$\times$ more packets in the 10\,ms
        window, and \texttt{HpHp\_L0.01\_radius} plus \texttt{HH\_L0.01\_pcc}
        carry 87\% of XGBoost's gain. Task-oriented devices are predictable, so
        the deviation is enormous rather than marginal.
  \item \textbf{Accuracy is the wrong headline metric for security.} The Isolation
        Forest's 84,437 false alerts/day and 7,562 missed attacks/day hide
        behind a single accuracy figure, and Logistic Regression's 0.9992
        accuracy conceals 4,553 false alerts/day. The error rates and their
        asymmetric costs are what an operator actually experiences, while any
        false negative leaves a device actively attacking from inside the
        network.
\end{enumerate}

Future work should add chronological splitting, adversarial robustness
testing, and a streaming evaluation harness to measure detection latency.

% =====================================================================
% \section{References}  -- rendered automatically by thebibliography below
% =====================================================================
\begin{thebibliography}{9}
% =====================================================================
\bibitem{kashif2020nbaiot}
M.~Kashif. \emph{N-BaIoT Dataset to Detect IoT Botnet Attacks} [Data set].
Kaggle, 2020. \url{https://www.kaggle.com/datasets/mkashifn/nbaiot-dataset}

\bibitem{meidan2018nbaiot}
Y.~Meidan, M.~Bohadana, Y.~Mathov, Y.~Mirsky, D.~Breitenbacher, A.~Shabtai,
and Y.~Elovici. N-BaIoT: Network-based Detection of IoT Botnet Attacks Using
Deep Autoencoders. \emph{arXiv:1805.03409}, 2018.

\bibitem{uci442}
UCI Machine Learning Repository. \emph{Detection of IoT Botnet Attacks N BaIoT}
(Dataset 442). \url{https://archive.ics.uci.edu/dataset/442}

\bibitem{kolias2017ddos}
C.~Kolias, G.~Kambourakis, A.~Stavrou, and J.~Voas. DDoS in the IoT: Mirai
and Other Botnets. \emph{Computer}, 50(7), 2017.

\bibitem{liu2008isolation}
F.~T.~Liu, K.~M.~Ting, and Z.-H.~Zhou. Isolation Forest. In \emph{Proc. IEEE
ICDM}, 2008.

\bibitem{breiman2001random}
L.~Breiman. Random Forests. \emph{Machine Learning}, 45(1):5--32, 2001.

\bibitem{pedregosa2011sklearn}
F.~Pedregosa et al.\ Scikit-learn: Machine Learning in Python. \emph{JMLR},
12:2825--2830, 2011.

\bibitem{chen2017xgboost}
T.~Chen and C.~Guestrin. XGBoost: A Scalable Tree Boosting System. In
\emph{Proc. ACM SIGKDD}, 2017.

\bibitem{lundberg2017shap}
S.~M.~Lundberg and S.-I.~Lee. A Unified Approach to Interpreting Model
Predictions. In \emph{NeurIPS}, 2017.
\end{thebibliography}

\end{document}
